\vspace{0.25cm}
Dans la partie « Diffusion de particules », l'accent est mis sur le notion de bilan dans le cas où le phénomène de convection est négligé. Cette partie se termine par un modèle microscopique.

{\tiny
\begin{minipage}[t]{0.5\textwidth}
    \begin{tabularx}{\textwidth}{|p{3cm}|X|}
    \hline 
    \multicolumn{1}{|c|}{Notions et contenus} & \multicolumn{1}{c|}{Capacités exigibles} \\
    \hline 
    3.2. Diffusion de particules & \\ 
    \hline 
    Vecteur densité de flux de particules. & Exprimer le flux de particules traversant une surface orientée en utilisant le vecteur $\Vec{\jN}$. \\
    \hline 
    Bilans de particules.  & Utiliser la notion de flux pour traduire un bilan global de particules. Établir l'équation locale traduisant un bilan de particules dans le cas d'un problème ne dépendant qu'une d'une seule coordonnée d'espace en coordonnées cartésiennes, cylindriques et sphériques, éventuellement en présence de sources internes. Utiliser l'opérateur divergence et son expression fournie pour exprimer le bilan local de particules dans le cas d'une géométrie quelconque. \\
    \hline 
    Loi de Fick. & Utiliser la loi de Fick. Citer l'ordre de grandeur d'un coefficient de diffusion dans un gaz dans les conditions usuelles. \\
    \hline 
    Régimes stationnaires. & Utiliser, en régime stationnaire, la conservation du flux sous forme locale ou globale en l'absence de sources internes. \\
    \hline 
    \end{tabularx}
\vfill
\end{minipage}%
\begin{minipage}[t]{0.5\textwidth}
    \begin{tabularx}{\textwidth}{|X|X|}
    \hline 
    \multicolumn{1}{|c|}{Notions et contenus} & \multicolumn{1}{c|}{Capacités exigibles} \\
   
    \hline 

    Équation de diffusion en l'absence de sources internes. & Établir l'équation de la diffusion en l'absence de sources internes. Utiliser l'opérateur laplacien et son expression fournie pour écrire l'équation de diffusion dans le cas d'une géométrie quelconque. Analyser une équation de diffusion en ordres de grandeur pour relier des échelles caractéristiques spatiale et temporelle. \\
    \hline 
    Approche microscopique du phénomène de diffusion. & 
Mettre en place un modèle probabiliste discret à une dimension de la diffusion (marche au hasard) et évaluer le coefficient de diffusion associé en fonction du libre parcours moyen et de la vitesse quadratique moyenne.
\textbf{Capacité numérique : à l'aide d'un langage de programmation, simuler la marche au hasard d'un grand nombre de particules à partir d'un centre et caractériser l'étalement spatial de cet ensemble de particules au cours du temps.} \\
\hline
\end{tabularx}
\vfill
\end{minipage}
}

